\fig{fv:n1:pulley}{弹簧秤承受的是两段绳张力的矢量和}{\begin{tikzpicture}[scale=.9]
\draw[thick](0,0)circle(.6);\draw[thick](-.6,0)--(-.6,-1.8) (.6,0)--(.6,-1.8);
\draw[->,blue,thick](-.6,-.4)--(-.6,-1.5)node[left]{$T$};\draw[->,blue,thick](.6,-.4)--(.6,-1.5)node[right]{$T$};
\draw[->,red,thick](0,.3)--(0,2)node[right]{$F_s$};\node at(3.6,.4){$F_s-2T=0$};\node at(3.6,-.5){轻滑轮的平动方程};
\end{tikzpicture}}
图\ref{fv:n1:pulley}只选滑轮为研究对象，两个向下的 $T$来自绳对轮的作用，向上的 $F_s$来自支架或秤。理想轻滑轮的质量趋零而受力有限，故合力必须趋零，秤读数为 $2T$。若两绳夹角改变，读数应改为两张力的合力大小；绳张力相等并不意味着支承读数永远是它们的代数和。另选绳端物体列方程时，不能再把 $F_s$当成该物体的受力。
